\chapter{شبکه‌های اسپایکی و محاسبات نورومورفیک}

\section{مبانی شبکه‌های عصبی اسپایکی}
پس از بحث سیستم پاداش و خطای پیش‌بینی پاداش، بحث وارد شبکه‌های عصبی اسپایکی یا \eng{Spiking Neural Networks} شد. هدف این بود که ابتدا معماری کلی این شبکه‌ها فهمیده شود و سپس ببینیم چگونه می‌توان سیستم پاداش را از دریچهٔ شبکه‌های اسپایکی مدل‌سازی کرد.

\subsection{چرا شبکه‌های اسپایکی؟}
در شبکه‌های عصبی مصنوعی معمول یا \eng{ANN}، ورودی‌ها اغلب اعداد پیوسته‌اند؛ برای مثال \eng{embedding}ها یا بردارهای عددی. این ورودی‌ها در وزن‌ها ضرب می‌شوند، جمع وزن‌دار ساخته می‌شود، \eng{bias} اضافه می‌شود، مقدار از \eng{activation function} عبور می‌کند و در نهایت خروجی یا پیش‌بینی تولید می‌شود. سپس پیش‌بینی با پاسخ درست یا \eng{ground truth} مقایسه می‌شود و اگر خطا وجود داشته باشد، وزن‌ها با \eng{backpropagation} اصلاح می‌شوند.

اما این تصویر از نظر زیستی کاملاً واقع‌گرایانه نیست. دو تفاوت مهم در این بحث برجسته شد: مغز با اعداد پیوسته‌ای شبیه ورودی‌های معمول \eng{ANN} کار نمی‌کند، و مغز \eng{backpropagation} کلاسیک و سراسری ندارد. در مغز، یادگیری بیشتر محلی است؛ هر نورون تا حد زیادی بر اساس ورودی‌ها و وضعیت خودش تصمیم می‌گیرد، نه اینکه کل شبکه منتظر خروجی نهایی بماند و سپس گرادیان سراسری به عقب برگردد.

شبکه‌های اسپایکی تلاش می‌کنند به این تصویر زیستی نزدیک‌تر شوند. در این شبکه‌ها، ارتباط نورون‌ها با رخدادهای گسستهٔ زمانی انجام می‌شود: اسپایک رخ می‌دهد یا رخ نمی‌دهد. به همین دلیل، این شبکه‌ها برای مدل‌سازی‌ای اهمیت دارند که می‌خواهد زمان، رویداد و مصرف انرژی را جدی بگیرد.

\subsection{تفاوت \eng{ANN} و \eng{SNN}}
در \eng{ANN} معمولاً با مقدارهای پیوسته سروکار داریم. نورون‌ها در هر گام مقدار عددی دریافت می‌کنند و خروجی عددی تولید می‌کنند. اما در \eng{SNN}، ورودی و خروجی نورون‌ها به‌شکل اسپایک است. اگر نورون اسپایک بزند، مقدار آن را می‌توان ۱ در نظر گرفت؛ اگر اسپایک نزند، مقدار ۰ است.

بنابراین در \eng{SNN}، ورودی یک دنبالهٔ زمانی از صفر و یک‌هاست:
\[
S = [0, 1, 0, 1, 1, 0, \ldots]
\]
وقتی مقدار ۱ باشد، یعنی اسپایک رخ داده است؛ و وقتی مقدار ۰ باشد، یعنی اسپایکی وجود ندارد. هر اسپایک در وزن سیناپسی ضرب می‌شود:
\[
S_i \times W_i
\]
اگر مجموع ورودی‌های وزن‌دار باعث شود ولتاژ نورون از آستانه عبور کند، نورون پس‌سیناپسی نیز اسپایک می‌زند.

\coursefigure{0.70}{artificial-neuron.png}{نورون مصنوعی کلاسیک در \eng{ANN}: ورودی‌های پیوستهٔ وزن‌دار جمع می‌شوند، از تابع فعال‌ساز عبور می‌کنند و خروجی عددی تولید می‌شود.}

\coursefigure{0.70}{spiking-neuron.png}{نورون اسپایکی در \eng{SNN}: ورودی‌های وزن‌دار به پتانسیل غشا تبدیل می‌شوند و خروجی به‌صورت رخداد شلیک اسپایک ظاهر می‌شود.}

\subsection{نورون زیستی}
برای فهم \eng{SNN}، ابتدا ساختار سادهٔ نورون زیستی مرور شد. یک نورون زیستی را می‌توان به‌صورت ساده شامل دندریت، بدنهٔ سلولی، آکسون و سیناپس دانست.

\subsubsection{دندریت}
دندریت محل دریافت ورودی است. سیگنال‌های نورون‌های دیگر از طریق سیناپس‌ها به دندریت می‌رسند.

\subsubsection{\eng{Soma}}
\eng{Soma} یا بدنهٔ سلولی بخش پردازشی نورون است. در توضیح ساده، اثر ورودی‌ها در این بخش جمع می‌شود و وضعیت نورون برای اسپایک زدن یا نزدن تعیین می‌شود.

\subsubsection{آکسون}
آکسون خروجی نورون را به نورون‌های دیگر منتقل می‌کند. اگر نورون اسپایک بزند، این رخداد از مسیر آکسون به سمت سیناپس‌ها منتقل می‌شود.

\subsubsection{سیناپس}
سیناپس محل اتصال آکسون یک نورون به دندریت نورون دیگر است. اگر نورونی به نورون دیگر اطلاعات بفرستد، نورون اول \eng{presynaptic neuron} و نورون دوم \eng{postsynaptic neuron} نامیده می‌شود. نورون پیش‌سیناپسی اطلاعات را می‌فرستد و نورون پس‌سیناپسی آن را دریافت می‌کند.

\subsection{اسپایک و زمان‌محور بودن محاسبه}
در \eng{SNN} زمان بخشی اصلی از محاسبه است. در \eng{ANN} معمولاً با یک بردار عددی برای یک شیء یا ورودی کار می‌کنیم، اما در \eng{SNN} ورودی یک توالی زمانی است. مهم است کدام اسپایک زودتر آمده، کدام بعداً آمده و فاصلهٔ زمانی میان آن‌ها چقدر بوده است.

به همین دلیل، خروجی \eng{SNN} نیز فقط یک مقدار عددی ساده نیست؛ بلکه معمولاً باید الگوی اسپایک‌ها در طول زمان تفسیر شود. هر اسپایک می‌تواند در همان لحظه بخشی از تصمیم‌گیری را فعال کند، بی‌آنکه لازم باشد همهٔ نورون‌های شبکه درگیر محاسبه شوند.

\subsection{مدل \eng{LIF}}
برای مدل‌سازی نورون اسپایکی، مدل‌های مختلفی وجود دارد؛ از جمله \eng{Leaky Integrate-and-Fire} یا \eng{LIF}، مدل \eng{Izhikevich} و مدل \eng{Hodgkin-Huxley}. در این بخش تمرکز روی مدل \eng{LIF} بود، چون هم ساده است و هم برخی رفتارهای مهم نورون زیستی را نشان می‌دهد.

فرمول سادهٔ مدل چنین است:
\[
V(t+1) = \lambda V(t) + I(t)
\]
در این رابطه، \(V(t)\) ولتاژ قبلی، \(\lambda\) ضریب نشت یا \eng{leak factor} و \(I(t)\) جریان ورودی است. جریان ورودی معمولاً از اسپایک‌های ورودی و وزن‌های سیناپسی ساخته می‌شود:
\[
I(t) = \sum_i W_i S_i(t)
\]
اگر ولتاژ از آستانه عبور کند:
\[
V(t+1) \geq V_{\text{\lr{threshold}}}
\]
نورون اسپایک می‌زند.

\coursefigure{0.78}{leaky-integrate-and-fire-neuron.png}{مدل \eng{LIF}: اسپایک‌های ورودی وزن‌دار در دندریت و سوما جمع می‌شوند، ولتاژ غشا با نشت زمانی تغییر می‌کند و پس از عبور از آستانه اسپایک تولید می‌شود.}

\subsubsection{\eng{Integration}}
\eng{Integration} یعنی اثر ورودی‌ها با هم جمع می‌شود. اگر نورون از چند نورون پیش‌سیناپسی ورودی بگیرد، اسپایک‌های ورودی در وزن‌های سیناپسی ضرب می‌شوند. سپس اثر آن‌ها در ولتاژ نورون تجمع پیدا می‌کند.

\subsubsection{\eng{Fire}}
\eng{Fire} یعنی اگر ولتاژ تجمع‌یافته از آستانه عبور کند، نورون اسپایک می‌زند. در این لحظه خروجی نورون برابر با رخداد اسپایک است.

\subsubsection{\eng{Leak}}
\eng{Leak} یعنی اگر ورودی کافی وجود نداشته باشد، ولتاژ نورون به‌تدریج کاهش پیدا می‌کند. این کاهش تدریجی همان نشت است و باعث می‌شود اثر ورودی‌ها برای همیشه در نورون باقی نماند.

\subsubsection{\eng{Refractory Period}}
پس از اسپایک، نورون برای مدت کوتاهی وارد دوره‌ای می‌شود که در آن فعالیتی ندارد یا به‌سادگی دوباره اسپایک نمی‌زند. این دوره \eng{refractory period} نامیده می‌شود و از رفتار نورون‌های واقعی الهام گرفته شده است.

\subsection{مصرف انرژی و پردازش رویدادمحور}
یکی از ویژگی‌های مهم \eng{SNN} این است که شبکه می‌تواند همیشه در حالت آماده‌باش یا \eng{on-call} باشد، اما مصرف انرژی پایینی داشته باشد. دلیل آن این است که وقتی ورودی صفر است، محاسبهٔ سنگین انجام نمی‌شود. نورون فقط در صورت وجود اسپایک و عبور از آستانه فعال می‌شود.

در \eng{ANN} معمولاً در هر گام محاسباتی، لایه‌ها درگیر ضرب و جمع می‌شوند، حتی اگر تغییر مهمی در ورودی رخ نداده باشد. اما در \eng{SNN} رخدادهای مهم، یعنی اسپایک‌ها، محاسبه را فعال می‌کنند. به همین دلیل، این شبکه‌ها برای سیستم‌هایی مناسب‌اند که باید همیشه فعال باشند اما مصرف انرژی کمی داشته باشند.

مثال‌های مناسب شامل دستیار صوتی همیشه‌روشن برای تشخیص \eng{wake-up word}، سمعک‌ها، سیستم‌های شنود کم‌مصرف، پهپادها، ربات‌هایی که باید نسبت به مانع یا حمله واکنش سریع نشان دهند، و سیستم‌هایی بود که به واکنش سریع و محلی نیاز دارند.

\section{یادگیری در شبکه‌های اسپایکی}
پس از معرفی معماری \eng{SNN}، بحث به این پرسش رسید که شبکهٔ اسپایکی چگونه یاد می‌گیرد. در مغز، وقتی برخی مسیرهای نورونی بارها فعال می‌شوند، اتصال میان آن‌ها تقویت می‌شود. بنابراین یادگیری در این سطح را می‌توان به تغییر وزن‌های سیناپسی ربط داد:
\[
\Delta W
\]
وقتی یک مسیر بارها استفاده شود، وزن سیناپسی آن افزایش پیدا می‌کند و آن مسیر در شبکه تقویت می‌شود.

برای توضیح این ایده، می‌توان از مثال مسیر روی چمن استفاده کرد. فرض کنید در یک دانشگاه مسیر رسمی برای عبور افراد طراحی شده است، اما پس از مدتی مسیر دیگری روی چمن بیشتر پا می‌خورد و سبزه‌هایش از بین می‌رود. این یعنی افراد بدون دستور مستقیم، مسیر دیگری را بارها انتخاب کرده‌اند. در شبکهٔ نورونی نیز مسیری که بارها فعال شود، وزن بیشتری می‌گیرد و تقویت می‌شود.

\subsection{قانون \eng{Hebbian Learning}}
ایدهٔ پایه برای یادگیری بدون‌ناظر در شبکه‌های اسپایکی، قانون \eng{Hebb} است. عبارت مشهور آن چنین است:
\[
\text{\lr{Fire together, wire together}}
\]
یعنی اگر دو نورون مرتب با هم فعال شوند، اتصال بین آن‌ها تقویت می‌شود.

در اینجا منظور از «با هم فعال شدن» الزاماً هم‌زمانی دقیق نیست. منظور این است که نورون پیش‌سیناپسی اسپایک بزند و نورون پس‌سیناپسی در ادامه فعال شود. اگر نورون پیش‌سیناپسی اسپایک بزند و بعد نورون پس‌سیناپسی نیز اسپایک بزند، وزن سیناپسی میان آن‌ها افزایش پیدا می‌کند.

این اتفاق مهم است، چون نشان می‌دهد میان آن دو نورون نوعی ارتباط اطلاعاتی وجود دارد. ممکن است هر دو نسبت به موضوعی مشابه حساس باشند، یا میان آن‌ها رابطه‌ای علّی وجود داشته باشد؛ یعنی فعال شدن نورون پیش‌سیناپسی در فعال شدن نورون پس‌سیناپسی نقش داشته باشد.

\subsection{\eng{STDP}}
\eng{STDP} مخفف \eng{Spike-Timing-Dependent Plasticity} است. یعنی تغییرپذیری وزن‌ها به زمان‌بندی اسپایک‌ها وابسته است. قانون \eng{Hebb} می‌گوید اگر دو نورون با هم فعال شوند، اتصالشان تقویت می‌شود. اما \eng{STDP} نکتهٔ دقیق‌تری اضافه می‌کند: ترتیب و فاصلهٔ زمانی اسپایک‌ها مهم است.

اگر نورون پیش‌سیناپسی پیش از نورون پس‌سیناپسی اسپایک بزند و به فعال شدن آن کمک کند، وزن سیناپسی باید تقویت شود. اما اگر نورون پس‌سیناپسی پیش از نورون پیش‌سیناپسی فعال شده باشد، احتمالاً فعال شدن پس‌سیناپسی به علت آن نورون پیش‌سیناپسی خاص نبوده است؛ بنابراین وزن میان آن دو باید تضعیف شود.

\subsubsection{\eng{LTP}}
\eng{LTP} یا \eng{Long-Term Potentiation} زمانی رخ می‌دهد که نورون پیش‌سیناپسی قبل از نورون پس‌سیناپسی اسپایک بزند:
\[
t_{\text{\lr{pre}}} < t_{\text{\lr{post}}}
\]
در این حالت رابطهٔ علّی-زمانی معقولی وجود دارد و وزن سیناپسی تقویت می‌شود:
\[
\Delta W > 0
\]

\subsubsection{\eng{LTD}}
\eng{LTD} یا \eng{Long-Term Depression} زمانی رخ می‌دهد که نورون پیش‌سیناپسی پس از نورون پس‌سیناپسی اسپایک بزند:
\[
t_{\text{\lr{pre}}} > t_{\text{\lr{post}}}
\]
در این حالت نورون پیش‌سیناپسی احتمالاً علت فعال شدن نورون پس‌سیناپسی نبوده است، پس وزن سیناپسی تضعیف می‌شود:
\[
\Delta W < 0
\]

\subsubsection{نقش اختلاف زمانی}
در \eng{STDP} فقط مثبت یا منفی بودن اختلاف زمانی مهم نیست؛ اندازهٔ اختلاف زمانی نیز اهمیت دارد. اختلاف زمانی را می‌توان چنین تعریف کرد:
\[
\Delta t = t_{\text{\lr{post}}} - t_{\text{\lr{pre}}}
\]
اگر \(\Delta t > 0\)، یعنی نورون پیش‌سیناپسی زودتر از نورون پس‌سیناپسی فعال شده و این حالت به \eng{LTP} مربوط است. اگر \(\Delta t < 0\)، یعنی نورون پس‌سیناپسی زودتر فعال شده و این حالت به \eng{LTD} مربوط است.

هرچه \(|\Delta t|\) کوچک‌تر باشد، دو اسپایک از نظر زمانی به هم نزدیک‌تر بوده‌اند و ارتباط میان آن‌ها قوی‌تر فرض می‌شود؛ پس تغییر وزن بیشتر خواهد بود. هرچه \(|\Delta t|\) بزرگ‌تر باشد، ارتباط ضعیف‌تر است و تغییر وزن کم‌تر می‌شود.

فرمول شهودی برای تقویت چنین بیان شد:
\[
\Delta W = A_+ e^{-\Delta t / \tau_+}
\]
و برای تضعیف:
\[
\Delta W = -A_- e^{\Delta t / \tau_-}
\]
در این فرمول‌ها، \(A_+\) دامنهٔ تقویت، \(A_-\) دامنهٔ تضعیف، \(\tau_+\) و \(\tau_-\) ثابت‌های زمانی، و \(\Delta t\) اختلاف زمانی اسپایک‌ها هستند.

اگر نورون پیش‌سیناپسی قبل از پس‌سیناپسی و با فاصلهٔ کم اسپایک بزند، تقویت زیاد است. اگر همین ترتیب با فاصلهٔ زیاد رخ دهد، تقویت کم‌تر است. اگر پیش‌سیناپسی بعد از پس‌سیناپسی و با فاصلهٔ کم اسپایک بزند، تضعیف زیاد است؛ و اگر فاصله زیاد باشد، تضعیف کم‌تر است.

اختلاف زمانی زمانی محاسبه می‌شود که هر دو نورون اسپایک زده باشند. اگر فقط نورون پیش‌سیناپسی اسپایک بزند و نورون پس‌سیناپسی فعال نشود، هنوز اختلاف زمانی معناداری برای \eng{STDP} وجود ندارد.

\subsection{\eng{Reward-Modulated STDP}}
تا اینجا دو جزء جداگانه داشتیم: از سیستم پاداش و \eng{TD Learning}، مفهوم خطای پیش‌بینی پاداش را داشتیم:
\[
\text{\lr{Reward Prediction Error}} =
R_{\text{\lr{actual}}} - R_{\text{\lr{predicted}}}
\]
و از شبکه‌های اسپایکی، تغییر وزن بر اساس زمان‌بندی اسپایک‌ها را داشتیم:
\[
\Delta W_{\text{\lr{STDP}}}
\]
هدف \eng{Reward-Modulated STDP} این است که سیستم پاداش روی \eng{STDP} سوار شود. یعنی هم اثر زمان‌بندی اسپایک‌ها لحاظ شود و هم اثر خطای پیش‌بینی پاداش.

فرم کلی تغییر وزن در این بحث چنین بیان شد:
\[
\Delta W = \eta \times \delta \times E(t)
\]
در این رابطه، \(\eta\) نرخ یادگیری، \(\delta\) خطای پیش‌بینی پاداش، \(E(t)\) اثر \eng{STDP} یا \eng{eligibility trace}، و \(\Delta W\) تغییر وزن سیناپسی است.

در این مدل، \eng{STDP} مشخص می‌کند کدام سیناپس‌ها از نظر زمانی در رفتار نقش داشته‌اند، و خطای پیش‌بینی پاداش مشخص می‌کند که نتیجهٔ رفتار باید آن سیناپس‌ها را تقویت کند یا تضعیف. اگر سیناپسی از نظر \eng{STDP} فعال و مرتبط بوده باشد و سپس سیستم پاداش بگوید نتیجه خوب بوده است، آن سیناپس بیشتر تقویت می‌شود. اما اگر نتیجه بدتر از انتظار باشد، همان سیناپس می‌تواند تضعیف شود.

این ایده همچنین به مسئلهٔ نسبت دادن علت‌ها کمک می‌کند. وقتی نورون پس‌سیناپسی اسپایک می‌زند، ممکن است از چند نورون پیش‌سیناپسی ورودی گرفته باشد و دقیقاً ندانیم کدام ورودی علت اصلی عبور از آستانه بوده است. \eng{STDP} ردپای نورون‌هایی را که کمی قبل‌تر اسپایک زده‌اند نگه می‌دارد، و پاداش بعدی تعیین می‌کند این ردپا در جهت تقویت استفاده شود یا تضعیف.

\subsection{از \eng{SNN} پایه به \eng{Deep SNN}}
در بخش‌های قبل، نورون اسپایکی با مدل \eng{LIF} و یادگیری سیناپسی با \eng{STDP} و \eng{Reward-Modulated STDP} بررسی شد. در این بخش، بحث از شبکه‌های کم‌عمق و پایه‌ای به سمت \eng{Deep SNN} رفت؛ یعنی شبکه‌هایی که چندین لایهٔ اسپایکی پشت سر هم دارند.

در یادآوری زیستی، نورون از دندریت، جسم سلولی، آکسون، پایانهٔ آکسونی، شکاف سیناپسی و وزن سیناپسی میان نورون پیش‌سیناپسی و پس‌سیناپسی تشکیل می‌شود. وقتی نورون پیش‌سیناپسی اسپایک می‌زند، این رویداد از راه وزن سیناپسی روی نورون بعدی اثر می‌گذارد و به شکل جریان ورودی در نورون پس‌سیناپسی ظاهر می‌شود.

در مدل \eng{LIF}، نورون جریان‌های ورودی را جمع می‌کند و ولتاژ غشایی آن افزایش می‌یابد. اگر ولتاژ از آستانه عبور کند، نورون اسپایک می‌زند:
\[
V > \theta
\]
پس از اسپایک، ولتاژ ریست می‌شود و نورون وارد دورهٔ \eng{refractory} می‌شود؛ یعنی برای مدتی دوباره فعال نمی‌شود یا فعال‌شدن آن دشوارتر است. خروجی سادهٔ نورون را می‌توان چنین نوشت:
\[
S =
\begin{cases}
1 & V > \theta \\
0 & V \leq \theta
\end{cases}
\]
بنابراین خروجی نورون اسپایکی معمولاً دودویی است: یا اسپایک داریم یا نداریم.

در ساده‌ترین حالت، اگر ورودی نورون \(x\) و وزن سیناپسی \(w\) باشد، جریان یا اثر سیناپسی چنین نوشته می‌شود:
\[
I = xw
\]
و در حالت چند ورودی:
\[
V = x_1w_1 + x_2w_2 + x_3w_3 + \cdots
\]
اگر ولتاژ از آستانه عبور کند، خروجی ۱ می‌شود؛ در غیر این صورت خروجی صفر باقی می‌ماند.

در \eng{Deep SNN} چندین لایهٔ نورون اسپایکی پشت سر هم قرار می‌گیرند. خروجی هر لایه، یعنی دنباله‌ای از اسپایک‌ها، به لایهٔ بعدی منتقل می‌شود و در آنجا دوباره وزن سیناپسی، جریان، ولتاژ و تصمیم اسپایک‌زدن یا نزدن محاسبه می‌شود. از این نظر، شبکه شبیه شبکه‌های عمیق معمولی چندلایه است؛ اما ورودی و خروجی هر لایه به شکل اسپایکی و دودویی است.

\subsection{الهام زیستی از ساختار چندلایهٔ مغز}
ایدهٔ چندلایه‌بودن شبکه‌ها از ساختار مغز نیز الهام می‌گیرد. در سیستم بینایی انسان، ناحیه‌هایی مانند \eng{V1}، \eng{V2}، \eng{V3}، \eng{V4}، \eng{MT/V5} و \eng{IT} در پردازش اطلاعات بینایی نقش دارند. در لایه‌های اولیه، ویژگی‌های ساده‌تری مانند لبه، خط، رنگ و جهت استخراج می‌شوند. هرچه پردازش جلوتر می‌رود، بازنمایی‌ها پیچیده‌تر و انتزاعی‌تر می‌شوند: از لبه و خط به شکل‌های پیچیده‌تر، از شکل‌ها به چهره، و در سطح‌های بالاتر به هویت فرد یا مفهوم شیء.

پس عمیق‌تر شدن شبکه با توانایی ساخت \eng{representation}های پیچیده‌تر ارتباط دارد. البته در مغز، پردازش همیشه یک مسیر خطی واحد نیست. خروجی ناحیه‌های اولیه می‌تواند وارد مسیرهای مختلف شود؛ برخی لایه‌ها میان چند وظیفه مشترک‌اند و سپس مسیرها تخصصی‌تر می‌شوند. این تصویر به ایدهٔ \eng{multi-task learning} شباهت دارد.

در سیستم بینایی، مسیر \eng{ventral} بیشتر با تشخیص «چیستی» شیء و مسیر \eng{dorsal} بیشتر با «کجایی»، حرکت و تعامل با شیء مرتبط دانسته می‌شود. اگر یک مسیر آسیب ببیند، ممکن است فرد شیء را ببیند اما درست با آن تعامل نکند؛ یا برعکس، بتواند با اشیا تعامل کند اما نتواند دقیق تشخیص دهد چه هستند.

این ساختار چندلایه فقط به بینایی محدود نیست. در سیستم لامسه نیز ناحیه‌هایی مانند \eng{S1} یا \eng{Primary Somatosensory Cortex} و \eng{S2} یا \eng{Secondary Somatosensory Cortex} وجود دارند. در \eng{S1} اطلاعات اولیه‌تری مانند لمس، درد و دما پردازش می‌شود و در \eng{S2} اطلاعات حسی ترکیب‌شده‌تر و انتزاعی‌تر شکل می‌گیرند.

\coursefigure{0.72}{visual-processing-hierarchy.png}{سلسله‌مراتب پردازش بینایی از نواحی اولیه مانند \eng{V1} تا مسیرهای بالاتر مانند \eng{V4}، \eng{MT/V5} و قشر گیجگاهی تحتانی.}

\coursefigure{0.78}{dorsal-and-ventral-visual-streams.png}{تقسیم پردازش بینایی به مسیر پشتی \eng{Where/How} و مسیر شکمی \eng{What} که برای فهم الهام زیستی شبکه‌های چندلایه مفید است.}

\subsection{مصرف انرژی و انگیزهٔ \eng{Deep SNN}}
یکی از انگیزه‌های مهم برای مطالعهٔ \eng{SNN}، مصرف انرژی پایین مغز است. مغز انسان با وجود پیچیدگی زیاد، حدود ۲۰ وات انرژی مصرف می‌کند؛ در حالی که مدل‌های عمیق معمولی روی \eng{GPU}ها و سامانه‌های بزرگ می‌توانند انرژی بسیار بیشتری مصرف کنند.

در این بحث دو دلیل برای این بهینگی برجسته شد. نخست، اسپایکی‌بودن است: در مغز، ورودی و خروجی نورون‌ها پیوسته نیست و بیشتر به شکل رویدادهای صفر و یک دیده می‌شود:
\[
0 \quad \text{یا} \quad 1
\]
این دودویی و رویدادمحور بودن می‌تواند مصرف انرژی را کاهش دهد و سطحی از انتزاع ایجاد کند.

دلیل دوم، محلی بودن محاسبات است. در مغز، هر نورون تا حد زیادی بر اساس اطلاعات محلی خود تصمیم می‌گیرد. این وضعیت با شبکه‌های عمیق معمولی متفاوت است که در آن‌ها خطا به‌صورت سراسری محاسبه می‌شود و با \eng{backpropagation} از انتها به ابتدا برمی‌گردد. محلی بودن نیز یکی از عواملی است که به کاهش مصرف انرژی کمک می‌کند.

سؤال اصلی این است که آیا می‌توان \eng{Deep SNN}هایی ساخت که هم توانایی شبکه‌های عمیق را داشته باشند و هم در فاز اجرا به مزیت انرژی و رویدادمحوری مغز نزدیک‌تر شوند. برای پاسخ به این سؤال، باید روش یادگیری در شبکه‌های اسپایکی عمیق را بررسی کرد.

\subsection{محدودیت‌های \eng{STDP} در شبکه‌های عمیق}
\eng{STDP} برای شبکه‌های کوچک یا کم‌عمق می‌تواند مفید باشد، به‌ویژه وقتی یادگیری بدون ناظر است، برچسب و پاداش نداریم، و شبکه قرار است از الگوهای تکرارشوندهٔ ورودی یاد بگیرد. در \eng{Reward-Modulated STDP} نیز همان ایدهٔ زمان‌بندی اسپایک با سیگنال پاداش ترکیب می‌شود تا پاداش تعیین کند ردپای سیناپسی باید تقویت شود یا تضعیف.

اما وقتی تعداد لایه‌ها زیاد می‌شود، \eng{STDP} به‌تنهایی کافی نیست. همگرایی سخت‌تر می‌شود، دقت پایین می‌آید، تنظیم وزن‌ها فقط با قوانین محلی دشوار می‌شود و شبکه برای \eng{fine-tune} کردن وزن‌ها به روش‌های دقیق‌تری نیاز پیدا می‌کند. بنابراین در \eng{Deep SNN}، مخصوصاً در وظایف نظارت‌شده، معمولاً به روش‌هایی شبیه \eng{backpropagation} نیاز داریم.

\subsection{\eng{Surrogate Gradient}}
در یادگیری نظارت‌شده، داده‌ها و برچسب‌ها مشخص‌اند. شبکه ورودی را می‌گیرد، خروجی تولید می‌کند و آن را با \eng{label} مقایسه می‌کند. سپس وزن‌ها اصلاح می‌شوند. در شبکه‌های عمیق معمولی، این اصلاح با \eng{backpropagation} انجام می‌شود.

مشکل اصلی در \eng{SNN} این است که خروجی نورون یک تابع پله‌ای است:
\[
S =
\begin{cases}
1 & V > \theta \\
0 & V \leq \theta
\end{cases}
\]
تابع پله‌ای تقریباً همه‌جا مشتق صفر دارد. بنابراین اگر بخواهیم \eng{backpropagation} انجام دهیم، مشتق لازم برای به‌روزرسانی وزن‌ها از مسیر خروجی اسپایک به دست نمی‌آید. پرسش اصلی این است: چگونه می‌توان در شبکهٔ اسپایکی \eng{backpropagation} انجام داد، وقتی تابع اسپایک مشتق‌پذیر نیست؟

\subsubsection{مسئلهٔ مشتق‌پذیری اسپایک}
راه‌حل رایج، استفاده از \eng{Surrogate Gradient} است. ایده این است که در فاز \eng{forward}، شبکه همچنان اسپایکی باقی بماند و خروجی نورون واقعاً صفر یا یک باشد؛ اما در فاز \eng{backward}، برای محاسبهٔ گرادیان، تابع پله‌ای با تابعی نرم‌تر و مشتق‌پذیر جایگزین شود.

\subsubsection{تابع جایگزین}
تابع جایگزین یا \eng{surrogate function} می‌تواند شکلی شبیه \eng{sigmoid} یا تابع نرم دیگر داشته باشد. نقش آن این نیست که رفتار واقعی نورون را در فاز پیش‌رو تغییر دهد؛ نقش آن فقط فراهم کردن مشتقی است که در \eng{backpropagation} بتوان از آن استفاده کرد.

\coursefigure{0.65}{surrogate-gradient-functions.png}{نمونه‌هایی از مشتق‌های جایگزین برای تابع شلیک؛ تقریب‌هایی مانند \eng{SuperSpike} و \eng{SLAYER} در اطراف آستانه گرادیان قابل استفاده فراهم می‌کنند.}

پس در این چارچوب:
\begin{itemize}
  \item در \eng{forward}، تابع واقعی اسپایک استفاده می‌شود؛
  \item در \eng{backward}، از تابع جایگزین مشتق‌پذیر برای محاسبهٔ گرادیان استفاده می‌شود.
\end{itemize}

\subsubsection{تفاوت با \eng{Activation Function} معمولی}
در شبکه‌های عصبی معمولی، \eng{activation function} در خود \eng{forward pass} استفاده می‌شود. یعنی خروجی هر لایه واقعاً از تابعی مانند \eng{sigmoid}، \eng{ReLU} یا \eng{tanh} عبور می‌کند.

اما در \eng{SNN} با \eng{surrogate gradient}، تابع جایگزین فقط برای \eng{backpropagation} استفاده می‌شود. در فاز \eng{forward} همچنان اسپایک واقعی داریم، ورودی و خروجی همچنان صفر و یک هستند، و تابع \eng{surrogate} فقط برای محاسبهٔ مشتق و اصلاح وزن‌ها وارد می‌شود. بنابراین \eng{surrogate function} همان \eng{activation function} معمولی نیست، چون نقش اصلی آن در فاز \eng{backward} است، نه در فاز \eng{forward}.

\subsection{\eng{Feedforward} در \eng{Deep SNN}}
در فاز \eng{feedforward}، اسپایک‌های لایهٔ قبلی به لایهٔ بعدی منتقل می‌شوند. اگر خروجی لایهٔ قبلی را \(S\) و وزن سیناپسی میان دو لایه را \(W\) بگیریم، جریان ورودی به نورون بعدی در ساده‌ترین حالت چنین است:
\[
I = SW
\]
در حالت زمانی، جریان می‌تواند اثر مقدار قبلی خود را نیز حفظ کند:
\[
I[t+1] = \alpha I[t] + S[t]W
\]
در این رابطه، \(I[t]\) جریان قبلی، \(\alpha\) ضریب کاهش یا \eng{decay} جریان، \(S[t]\) اسپایک ورودی و \(W\) وزن سیناپسی است.

ولتاژ نورون نیز با توجه به ولتاژ قبلی، ضریب نشت و جریان ورودی به‌روزرسانی می‌شود:
\[
V[t+1] = \beta V[t] + I[t+1]
\]
که در آن، \(V[t]\) ولتاژ قبلی، \(\beta\) ضریب نشت ولتاژ و \(I[t+1]\) جریان ورودی جدید است.

اگر ولتاژ از آستانه بیشتر شود، اسپایک تولید می‌شود:
\[
S[t+1] =
\begin{cases}
1 & V[t+1] > \theta \\
0 & V[t+1] \leq \theta
\end{cases}
\]
پس از اسپایک، نورون باید ریست شود؛ یعنی ولتاژ و گاهی جریان آن کاهش یابد یا صفر شود.

\subsubsection{نقش \eng{reset}}
وقتی نورون اسپایک می‌زند، فقط به لایهٔ بعدی سیگنال نمی‌فرستد؛ بلکه خودش نیز ریست می‌شود. بنابراین اسپایک دو اثر دارد: به لایهٔ بعدی منتقل می‌شود و در وزن سیناپسی ضرب می‌شود، و هم‌زمان روی نورونی که اسپایک زده اثر می‌گذارد و ولتاژ آن را ریست می‌کند. به همین دلیل، در \eng{SNN} زمان و وضعیت داخلی نورون اهمیت اساسی دارند.

\subsection{\eng{Backpropagation} با \eng{Surrogate Gradient}}
در انتهای شبکه، خروجی با هدف مقایسه می‌شود و خطا محاسبه می‌شود. برای نمونه، اگر از \eng{Mean Squared Error} استفاده کنیم:
\[
E = \frac{1}{2}(T - O)^2
\]
که در آن \(T\) هدف و \(O\) خروجی شبکه است.

برای به‌روزرسانی وزن‌ها، مانند شبکه‌های معمولی از گرادیان استفاده می‌کنیم:
\[
W_{\text{\lr{new}}} =
W_{\text{\lr{old}}} - \eta \frac{\partial E}{\partial W}
\]
در این رابطه، \(\eta\) نرخ یادگیری و \(\frac{\partial E}{\partial W}\) مشتق خطا نسبت به وزن است.

با استفاده از قاعدهٔ زنجیره‌ای:
\[
\frac{\partial E}{\partial W}
=
\frac{\partial E}{\partial O}
\cdot
\frac{\partial O}{\partial Net}
\cdot
\frac{\partial Net}{\partial W}
\]
اگر:
\[
Net = XW
\]
آنگاه:
\[
\frac{\partial Net}{\partial W} = X
\]
و برای خطای مربعی:
\[
\frac{\partial E}{\partial O} = O - T
\]
بخش اصلی \eng{surrogate gradient} در جملهٔ میانی است:
\[
\frac{\partial O}{\partial Net}
\]
زیرا اگر \(O\) را خروجی تابع پله‌ای واقعی بگیریم، مشتق آن صفر می‌شود. بنابراین این مشتق را از تابع \eng{surrogate} می‌گیریم.

\subsubsection{مثال عددی}
فرض کنید:
\[
X = 2,\qquad W = 0.4,\qquad \theta = 1,\qquad T = 1,\qquad V_0 = 0
\]
ابتدا مقدار \(Net\) را حساب می‌کنیم:
\[
Net = XW = 2 \times 0.4 = 0.8
\]
چون:
\[
0.8 < 1
\]
نورون اسپایک نمی‌زند و بنابراین:
\[
O = 0
\]
اما هدف برابر ۱ بوده است:
\[
T = 1
\]
پس خروجی شبکه اشتباه است و وزن باید افزایش یابد.

مشتق خطا نسبت به خروجی:
\[
\frac{\partial E}{\partial O} = O - T = 0 - 1 = -1
\]
فرض شد مقدار مشتق تابع \eng{surrogate} در نقطهٔ \(Net=0.8\) تقریباً \(0.25\) است. همچنین:
\[
\frac{\partial Net}{\partial W} = X = 2
\]
پس:
\[
\frac{\partial E}{\partial W}
=
(-1)(0.25)(2)
=
-0.5
\]
اگر نرخ یادگیری \(\eta = 0.1\) باشد:
\[
W_{\text{\lr{new}}}
=
0.4 - 0.1(-0.5)
=
0.45
\]
بنابراین وزن از \(0.4\) به \(0.45\) افزایش پیدا می‌کند. این نتیجه منطقی است، چون خروجی باید ۱ می‌شد اما نورون به آستانه نرسیده بود؛ در نتیجه وزن افزایش می‌یابد تا دفعهٔ بعد ولتاژ راحت‌تر از آستانه عبور کند.

\subsection{محلی بودن، آموزش سراسری و ارزش عملی \eng{SNN}}
استفاده از \eng{surrogate gradient} در مرحلهٔ آموزش تا حدی ویژگی محلی بودن \eng{SNN} را کاهش می‌دهد، چون خطا به‌صورت سراسری محاسبه می‌شود و \eng{backpropagation} انجام می‌گیرد. اما نکتهٔ مهم این است که \eng{surrogate gradient} فقط در مرحلهٔ آموزش استفاده می‌شود.

پس از آموزش، شبکه همچنان اسپایکی است: ورودی و خروجی صفر و یک باقی می‌مانند، مدل می‌تواند روی سخت‌افزار \eng{neuromorphic} اجرا شود، و دیگر نیازی به \eng{surrogate gradient} یا \eng{backpropagation} در فاز اجرا نیست. وزن‌ها پیش‌تر \eng{fine-tune} شده‌اند.

بنابراین چارچوب اصلی \eng{SNN} حفظ می‌شود. فقط برای اینکه بتوانیم وزن‌ها را در شبکهٔ عمیق و نظارت‌شده تنظیم کنیم، در مرحلهٔ آموزش از \eng{surrogate gradient} کمک می‌گیریم.

ارزش عملی \eng{SNN} بیشتر در فاز اجرا روشن می‌شود. در کاربردهایی مانند ربات، پهپاد، سامانه‌های \eng{embedded}، سخت‌افزار \eng{neuromorphic} و پردازش \eng{real-time} با مصرف انرژی کم، شبکهٔ آموزش‌دیده می‌تواند به‌صورت اسپایکی اجرا شود. در این حالت، مهم نیست وزن‌ها دقیقاً با چه روشی آموزش دیده‌اند؛ مهم این است که \eng{inference} با اسپایک‌های صفر و یک انجام شود.

\subsection{کاربرد \eng{STDP} در برابر \eng{Surrogate Gradient}}
\eng{STDP} مناسب‌تر است وقتی شبکه کوچک یا کم‌عمق است، یادگیری بدون ناظر است، یادگیری \eng{online} لازم داریم، یا هدف، پیاده‌سازی زیستی‌تر است. در مقابل، \eng{surrogate gradient} برای \eng{Deep SNN} و وظایف نظارت‌شده مناسب‌تر است، جایی که باید خطای خروجی با هدف مقایسه شود و وزن‌ها در چندین لایه تنظیم شوند.
